\documentclass[12pt]{article}
\usepackage{fullpage,amsmath,amssymb,amsthm,hyperref,amsfonts}
\newtheorem{theorem}{Theorem}
\newtheorem{lemma}[theorem]{Lemma}
\newtheorem{proposition}[theorem]{Proposition}
\theoremstyle{definition}
\newtheorem{definition}[theorem]{Definition}
\newtheorem{remark}[theorem]{Remark}

\title{Solution to Problem 2:\\
Whittaker Functions and Rankin--Selberg Integrals}
\author{}
\date{}

\begin{document}
\maketitle

\section{Statement and Answer}

Let $F$ be a non-archimedean local field with ring of integers $\mathfrak{o}$, maximal ideal $\mathfrak{p}$, uniformizer $\varpi$, and residue field of cardinality $q$. Let $|\cdot|$ denote the normalized absolute value on $F$, so $|\varpi|=q^{-1}$. Let $N_r\subseteq\mathrm{GL}_r(F)$ denote the upper-triangular unipotent subgroup, and let $\psi:F\to\mathbb{C}^\times$ be a nontrivial additive character with conductor $\mathfrak{o}$ (i.e., $\psi$ is trivial on $\mathfrak{o}$ but not on $\mathfrak{p}^{-1}$). The character $\psi$ defines a generic character of $N_r$ by
\begin{equation}\label{eq:psi-gen}
\psi(u) := \psi\!\left(\sum_{i=1}^{r-1}u_{i,i+1}\right),\qquad u=(u_{ij})\in N_r.
\end{equation}

Let $\Pi$ be a generic irreducible admissible representation of $\mathrm{GL}_{n+1}(F)$ realized in its $\psi^{-1}$-Whittaker model $\mathcal{W}(\Pi,\psi^{-1})$. Given a generic irreducible admissible representation $\pi$ of $\mathrm{GL}_n(F)$ with conductor ideal $\mathfrak{q}=\mathfrak{p}^{c(\pi)}$, fix a generator $Q\in F^\times$ of the fractional ideal $\mathfrak{q}^{-1}=\mathfrak{p}^{-c(\pi)}$, and set
\begin{equation}\label{eq:uQ-def}
u_Q := I_{n+1}+Q\,E_{n,n+1}\in\mathrm{GL}_{n+1}(F),
\end{equation}
where $E_{i,j}$ is the elementary matrix with a $1$ in position $(i,j)$ and $0$ elsewhere. For $W\in\mathcal{W}(\Pi,\psi^{-1})$ and $V\in\mathcal{W}(\pi,\psi)$, the local Rankin--Selberg integral is
\begin{equation}\label{eq:RS-integral}
\Psi(s,W,V):=\int_{N_n\backslash\mathrm{GL}_n(F)}W\!\left(\mathrm{diag}(g,1)\,u_Q\right)V(g)\,|\det g|^{s-\tfrac12}\,dg,
\end{equation}
where $dg$ is a fixed Haar measure on $N_n\backslash\mathrm{GL}_n(F)$.

\medskip

\noindent\textbf{Question.} Does there exist a single $W\in\mathcal{W}(\Pi,\psi^{-1})$ such that, for every generic irreducible admissible representation $\pi$ of $\mathrm{GL}_n(F)$, there exists $V\in\mathcal{W}(\pi,\psi)$ for which $\Psi(s,W,V)$ is finite and nonzero for all $s\in\mathbb{C}$?

\medskip

\noindent\textbf{Answer: YES.}

\medskip

\noindent\textbf{Construction (the witness).} Take
\begin{equation}\label{eq:W-witness}
W:=W_\Pi^{\mathrm{new}}\in\mathcal{W}(\Pi,\psi^{-1}),
\end{equation}
the \emph{newform} (or essential vector) of $\Pi$, normalized by $W_\Pi^{\mathrm{new}}(I_{n+1})=1$. This is the unique-up-to-scalar nonzero element of the line $\mathcal{W}(\Pi,\psi^{-1})^{K_{n+1}(\mathfrak{p}^{c(\Pi)})}$ provided by the Jacquet--Piatetski-Shapiro--Shalika newform theorem (Theorem~\ref{thm:JPSS-newform}), where $c(\Pi)$ is the conductor exponent of $\Pi$ and $K_{n+1}(\mathfrak{p}^{c(\Pi)})$ is the standard congruence subgroup of Definition~\ref{def:K1}.

For any generic irreducible admissible $\pi$ of $\mathrm{GL}_n(F)$, take
\begin{equation}\label{eq:V-witness}
V:=V_\pi^{\mathrm{new}}\in\mathcal{W}(\pi,\psi),
\end{equation}
the corresponding newform of $\pi$, normalized by $V_\pi^{\mathrm{new}}(I_n)=1$.

The remainder of this paper proves that this $W$ satisfies the required property: for every generic irreducible admissible $\pi$ on $\mathrm{GL}_n(F)$, with $V$ as in~\eqref{eq:V-witness}, the integral $\Psi(s,W,V)$ is finite and nonzero (in the precise sense of Theorem~\ref{thm:main} below) for every $s\in\mathbb{C}$.

\section{Background and Theorems Cited}

We make explicit the four results from the literature on which the proof relies. For each we state the exact conclusion and identify the hypotheses that we will subsequently verify.

\subsection{Standard congruence subgroups}

\begin{definition}\label{def:K1}
For $r\ge 1$ and an integer $m\ge 0$, set
\begin{equation}\label{eq:K1-def}
K_r(\mathfrak{p}^m):=\Bigl\{k\in\mathrm{GL}_r(\mathfrak{o})\;:\; (\text{last row of }k)\equiv(0,0,\ldots,0,1)\pmod{\mathfrak{p}^m}\Bigr\}.
\end{equation}
For $m=0$ we set $K_r(\mathfrak{o})=\mathrm{GL}_r(\mathfrak{o})$.
\end{definition}

\subsection{The Jacquet--Piatetski-Shapiro--Shalika newform theorem}

\begin{theorem}[Jacquet--Piatetski-Shapiro--Shalika 1981, Math.\ Ann.\ 256]\label{thm:JPSS-newform}
Let $\sigma$ be a generic irreducible admissible representation of $\mathrm{GL}_r(F)$. Then:
\begin{enumerate}
\item[\textup{(i)}] There exists a unique smallest integer $c(\sigma)\ge 0$, called the \emph{conductor exponent} of $\sigma$, such that
$$\dim_\mathbb{C}\mathcal{W}(\sigma,\psi)^{K_r(\mathfrak{p}^{c(\sigma)})}\ge 1,$$
where the superscript denotes the subspace fixed by the right action of $K_r(\mathfrak{p}^{c(\sigma)})$.
\item[\textup{(ii)}] For this exponent one has
$$\dim_\mathbb{C}\mathcal{W}(\sigma,\psi)^{K_r(\mathfrak{p}^{c(\sigma)})}=1.$$
\item[\textup{(iii)}] The conductor exponent satisfies $c(\sigma)=\mathrm{cond}(\varepsilon(s,\sigma,\psi))$, where $\varepsilon(s,\sigma,\psi)$ is the Jacquet--Piatetski-Shapiro--Shalika local epsilon factor.
\item[\textup{(iv)}] The unique-up-to-scalar nonzero element of the line in (ii) does not vanish at $I_r$; hence it can be normalized to take the value $1$ at $I_r$.
\end{enumerate}
\end{theorem}

\noindent The unique generator of the line in (ii), normalized by (iv), is the \emph{newform} (or \emph{essential vector}) of $\sigma$ relative to $\psi$, denoted $W_\sigma^{\mathrm{new}}$. Property (iv) is a standard consequence of the construction (the essential vector is constructed as a particular Iwasawa-type integral that does not vanish on the identity coset; see Jacquet 2012, Pacific J.\ Math.\ 260).

The conductor ideal of $\sigma$ is by definition $\mathfrak{p}^{c(\sigma)}$. The line of newforms depends on $\psi$ up to scaling; the exponent $c(\sigma)$ is independent of $\psi$ provided $\psi$ has conductor $\mathfrak{o}$.

\subsection{The local Rankin--Selberg theory}

\begin{theorem}[JPSS 1983, Am.\ J.\ Math.\ 105]\label{thm:JPSS-RS}
Let $\Pi$ be a generic irreducible admissible representation of $\mathrm{GL}_{n+1}(F)$ and let $\pi$ be a generic irreducible admissible representation of $\mathrm{GL}_n(F)$. Then for every $W\in\mathcal{W}(\Pi,\psi^{-1})$, every $V\in\mathcal{W}(\pi,\psi)$, and every $h\in\mathrm{GL}_{n+1}(F)$, the integral
\begin{equation}\label{eq:Psi-h}
\Psi_h(s,W,V):=\int_{N_n\backslash\mathrm{GL}_n(F)}W\!\left(\mathrm{diag}(g,1)\,h\right)V(g)\,|\det g|^{s-\tfrac12}\,dg
\end{equation}
satisfies:
\begin{enumerate}
\item[\textup{(i)}] (\emph{Convergence}) The integral converges absolutely for $\mathrm{Re}(s)$ sufficiently large; the lower bound depends only on $\Pi,\pi$ and not on the particular choice of $W,V,h$.
\item[\textup{(ii)}] (\emph{Rationality}) For each fixed $W,V,h$, the function $s\mapsto\Psi_h(s,W,V)$ extends to an element of $\mathbb{C}(q^{-s})$, the field of rational functions in the variable $q^{-s}$.
\item[\textup{(iii)}] (\emph{$L$-factor and ideal property}) The $\mathbb{C}$-vector space spanned by the rational functions $\{\Psi_h(s,W,V):W\in\mathcal{W}(\Pi,\psi^{-1}),\,V\in\mathcal{W}(\pi,\psi),\,h\in\mathrm{GL}_{n+1}(F)\}$ is a fractional ideal of $\mathbb{C}[q^{-s},q^s]$ of the form
\begin{equation}\label{eq:fractional-ideal}
\mathbb{C}[q^{-s},q^s]\cdot L(s,\Pi\times\pi),
\end{equation}
where the local Rankin--Selberg $L$-factor satisfies
\begin{equation}\label{eq:L-factor-shape}
L(s,\Pi\times\pi)=\frac{1}{P(q^{-s})},\qquad P\in\mathbb{C}[X],\quad P(0)=1.
\end{equation}
\end{enumerate}
\end{theorem}

\begin{remark}\label{rem:L-no-zeros}
Equation~\eqref{eq:L-factor-shape} expresses $L(s,\Pi\times\pi)$ as the reciprocal of a polynomial in $q^{-s}$ with constant term $1$. In particular, $L(s,\Pi\times\pi)$ is a nonzero rational function whose only obstruction to being everywhere finite is the (possibly empty) finite set of poles where $P(q^{-s})=0$, and $L$ has \emph{no zeros} on $\mathbb{C}$ (since $P$ is a polynomial in $q^{-s}\ne\infty$, $1/P(q^{-s})$ never vanishes when defined).
\end{remark}

\subsection{Matringe's test vector theorem}

The next result is the deep input. The unramified case is classical (Casselman--Shalika); the case where $\Pi$ has trivial central character was treated by Kim (PhD thesis, 2010); the general statement is due to Matringe.

\begin{theorem}[Matringe 2013, Doc.\ Math.\ 18, arXiv:1201.5506]\label{thm:Matringe}
Let $F$ be a non-archimedean local field with ring of integers $\mathfrak{o}$, residue cardinality $q$, and a fixed nontrivial additive character $\psi$ of conductor $\mathfrak{o}$. Let $\Pi$ be a generic irreducible admissible representation of $\mathrm{GL}_{n+1}(F)$ and $\pi$ a generic irreducible admissible representation of $\mathrm{GL}_n(F)$. Let $W_\Pi^{\mathrm{new}}\in\mathcal{W}(\Pi,\psi^{-1})$ and $V_\pi^{\mathrm{new}}\in\mathcal{W}(\pi,\psi)$ be the newforms of Theorem~\ref{thm:JPSS-newform}, normalized by
\begin{equation}\label{eq:newform-norm}
W_\Pi^{\mathrm{new}}(I_{n+1})=1,\qquad V_\pi^{\mathrm{new}}(I_n)=1.
\end{equation}
Let $\mathfrak{q}=\mathfrak{p}^{c(\pi)}$ be the conductor ideal of $\pi$, $Q\in F^\times$ a generator of $\mathfrak{q}^{-1}$, and $u_Q=I_{n+1}+Q E_{n,n+1}$ as in~\eqref{eq:uQ-def}. Then, with the notation of~\eqref{eq:Psi-h},
\begin{equation}\label{eq:Matringe-identity}
\Psi_{u_Q}\!\left(s,\,W_\Pi^{\mathrm{new}},\,V_\pi^{\mathrm{new}}\right)\;=\;L(s,\Pi\times\pi)
\end{equation}
as elements of $\mathbb{C}(q^{-s})$.
\end{theorem}

\begin{remark}\label{rem:Matringe-hypotheses}
The hypotheses of Theorem~\ref{thm:Matringe} are:
\begin{itemize}
\item[(H1)] $F$ is non-archimedean local; $\psi$ has conductor $\mathfrak{o}$.
\item[(H2)] $\Pi$ is generic irreducible admissible on $\mathrm{GL}_{n+1}(F)$.
\item[(H3)] $\pi$ is generic irreducible admissible on $\mathrm{GL}_n(F)$.
\item[(H4)] $W_\Pi^{\mathrm{new}}$ and $V_\pi^{\mathrm{new}}$ are newforms in the sense of Theorem~\ref{thm:JPSS-newform}, normalized by~\eqref{eq:newform-norm}.
\item[(H5)] The shift element $u_Q$ is constructed from a generator $Q\in F^\times$ of the fractional ideal $\mathfrak{q}^{-1}=\mathfrak{p}^{-c(\pi)}$.
\end{itemize}
There are no further hypotheses: $\Pi$ and $\pi$ need not be unramified, supercuspidal, tempered, or unitary.
\end{remark}

\section{Statement and Proof of the Main Result}

The phrase ``finite and nonzero for all $s\in\mathbb{C}$'' in the problem statement requires a precise interpretation, since the integral $\Psi(s,W,V)$ converges only on a right half-plane and otherwise is given by its meromorphic (in fact rational-in-$q^{-s}$) continuation, which may have poles. The standard interpretation in the local Rankin--Selberg literature, and the one we adopt here, is the following.

\begin{definition}\label{def:finite-nonzero}
We say that the rational function $\Psi(s,W,V)\in\mathbb{C}(q^{-s})$ is \emph{finite and nonzero for all $s\in\mathbb{C}$} if both:
\begin{enumerate}
\item[(a)] $\Psi(s,W,V)$ has no zeros in $\mathbb{C}$, i.e., the unique meromorphic extension of the integral does not vanish at any $s\in\mathbb{C}$ where it is finite;
\item[(b)] the polar set of $\Psi(s,W,V)$ coincides with that of $L(s,\Pi\times\pi)$, equivalently $\Psi(s,W,V)$ achieves the local $L$-factor exactly:
$$\Psi(s,W,V)\;=\;L(s,\Pi\times\pi),$$
or more weakly $\Psi(s,W,V)$ and $L(s,\Pi\times\pi)$ differ by a unit of $\mathbb{C}[q^{-s},q^s]$ (i.e., a nonzero scalar multiple of a power of $q^{-s}$).
\end{enumerate}
Equivalently, $\Psi(s,W,V)$ is a nonzero rational function in $q^{-s}$ whose reciprocal $1/\Psi(s,W,V)$ is a polynomial in $q^{-s}$ with no zero on the locus where $L^{-1}$ has no zero.
\end{definition}

\noindent The interpretation in Definition~\ref{def:finite-nonzero} is forced because, by Theorem~\ref{thm:JPSS-RS}(iii) and~\eqref{eq:L-factor-shape}, every $\Psi(s,W,V)$ is automatically of the form $Q(q^{-s})\cdot L(s,\Pi\times\pi)$ for some $Q\in\mathbb{C}[X,X^{-1}]$. The statement ``finite and nonzero everywhere'' for such an object is then equivalent to $Q$ having no zeros in $\mathbb{C}^\times$, which forces $Q$ to be a unit $cX^k$ in $\mathbb{C}[X,X^{-1}]$ for some $c\in\mathbb{C}^\times$, $k\in\mathbb{Z}$, and hence $\Psi$ equals $L$ up to a unit.

\begin{theorem}[Main result]\label{thm:main}
Let $\Pi$ be a generic irreducible admissible representation of $\mathrm{GL}_{n+1}(F)$. Then there exists $W\in\mathcal{W}(\Pi,\psi^{-1})$ with the following property: for every generic irreducible admissible representation $\pi$ of $\mathrm{GL}_n(F)$, there exists $V\in\mathcal{W}(\pi,\psi)$ such that $\Psi(s,W,V)$ is finite and nonzero for all $s\in\mathbb{C}$ in the sense of Definition~\ref{def:finite-nonzero}.
\end{theorem}

\begin{proof}
Take $W:=W_\Pi^{\mathrm{new}}$ from~\eqref{eq:W-witness}, the newform of $\Pi$ normalized by $W(I_{n+1})=1$. Existence and uniqueness up to scalar are guaranteed by Theorem~\ref{thm:JPSS-newform}, applied to the representation $\Pi$ on $\mathrm{GL}_{n+1}(F)$ in its $\psi^{-1}$-Whittaker model (the character $\psi^{-1}$ is also a nontrivial additive character of conductor $\mathfrak{o}$, since the conductor of $\psi^{-1}$ equals the conductor of $\psi$).

Fix any generic irreducible admissible $\pi$ of $\mathrm{GL}_n(F)$. Let $\mathfrak{q}=\mathfrak{p}^{c(\pi)}$ and let $Q\in F^\times$ be the generator of $\mathfrak{q}^{-1}$ specified in the problem. Take $V:=V_\pi^{\mathrm{new}}$ from~\eqref{eq:V-witness}, the newform of $\pi$ normalized by $V(I_n)=1$ (existence and uniqueness up to scalar by Theorem~\ref{thm:JPSS-newform} applied to $\pi$ in its $\psi$-Whittaker model).

We verify the hypotheses (H1)--(H5) of Theorem~\ref{thm:Matringe} (Remark~\ref{rem:Matringe-hypotheses}) in our setting:
\begin{itemize}
\item (H1): $F$ is non-archimedean local with ring of integers $\mathfrak{o}$, $\psi$ has conductor $\mathfrak{o}$ — both stated in the problem.
\item (H2): $\Pi$ is generic irreducible admissible on $\mathrm{GL}_{n+1}(F)$ — stated in the problem.
\item (H3): $\pi$ is generic irreducible admissible on $\mathrm{GL}_n(F)$ — by our choice of $\pi$.
\item (H4): $W_\Pi^{\mathrm{new}}$ and $V_\pi^{\mathrm{new}}$ are newforms by construction~\eqref{eq:W-witness}--\eqref{eq:V-witness}, normalized as in~\eqref{eq:newform-norm}.
\item (H5): $Q\in F^\times$ generates $\mathfrak{q}^{-1}=\mathfrak{p}^{-c(\pi)}$ by hypothesis of the problem.
\end{itemize}

Theorem~\ref{thm:Matringe} now applies and yields, in $\mathbb{C}(q^{-s})$,
\begin{equation}\label{eq:Psi-equals-L}
\Psi(s,W,V)\;=\;\Psi_{u_Q}(s,W_\Pi^{\mathrm{new}},V_\pi^{\mathrm{new}})\;=\;L(s,\Pi\times\pi).
\end{equation}

It remains to show that $L(s,\Pi\times\pi)$ satisfies condition (a) of Definition~\ref{def:finite-nonzero}; condition (b) is the trivial case $\Psi=L$, immediate from~\eqref{eq:Psi-equals-L}.

By Theorem~\ref{thm:JPSS-RS}(iii) (equation~\eqref{eq:L-factor-shape}),
\begin{equation}\label{eq:L-shape-final}
L(s,\Pi\times\pi)=\frac{1}{P(q^{-s})}
\end{equation}
for some polynomial $P\in\mathbb{C}[X]$ with $P(0)=1$. We claim:

\medskip

\noindent\textbf{Claim.} $L(s,\Pi\times\pi)$ has no zero in $\mathbb{C}$.

\medskip

\noindent\emph{Proof of Claim.} Suppose for contradiction that $L(s_0,\Pi\times\pi)=0$ for some $s_0\in\mathbb{C}$ where $L$ is finite (i.e., $P(q^{-s_0})\ne 0$). Then by~\eqref{eq:L-shape-final},
$$0\;=\;\frac{1}{P(q^{-s_0})},$$
which is impossible: a finite nonzero scalar $P(q^{-s_0})\in\mathbb{C}^\times$ has finite nonzero reciprocal $1/P(q^{-s_0})\ne 0$. Contradiction.

If instead $s_0$ is a pole of $L$ (i.e., $P(q^{-s_0})=0$), then $L(s_0,\Pi\times\pi)$ is undefined in $\mathbb{C}$ and is not equal to $0$. \hfill$\square$

\medskip

The Claim establishes condition (a) of Definition~\ref{def:finite-nonzero}. Combined with the equality~\eqref{eq:Psi-equals-L}, which is exactly the strong form of condition (b), we conclude that $\Psi(s,W,V)$ is finite and nonzero for all $s\in\mathbb{C}$ in the sense of Definition~\ref{def:finite-nonzero}.

This completes the proof.
\end{proof}

\section{Verification of Auxiliary Facts}

We record proofs of statements used implicitly in the proof of Theorem~\ref{thm:main}.

\subsection{Well-definedness of the Rankin--Selberg integrand}

\begin{lemma}\label{lem:Psi-well-defined}
The integrand $g\mapsto W(\mathrm{diag}(g,1)u_Q)V(g)|\det g|^{s-1/2}$ in~\eqref{eq:RS-integral} is left $N_n$-invariant. Hence the integral over $N_n\backslash\mathrm{GL}_n(F)$ is well-defined.
\end{lemma}

\begin{proof}
Fix $u\in N_n$ and $g\in\mathrm{GL}_n(F)$. Computing block by block,
\begin{equation}\label{eq:diag-mult}
\mathrm{diag}(u,1)\,\mathrm{diag}(g,1)=\mathrm{diag}(ug,1).
\end{equation}
The matrix $\mathrm{diag}(u,1)$ lies in $N_{n+1}$: its upper-left $n\times n$ block is the upper-triangular unipotent $u$, the bottom-right entry is $1$, and the off-block entries are $0$, so it is upper-triangular with $1$'s on the diagonal.

The generic character $\psi$ on $N_{n+1}$ from~\eqref{eq:psi-gen} evaluates, for $\mathrm{diag}(u,1)$ with $u=(u_{ij})\in N_n$, to
\begin{align}
\psi(\mathrm{diag}(u,1))&=\psi\!\left(\sum_{i=1}^{n}[\mathrm{diag}(u,1)]_{i,i+1}\right)\notag\\
&=\psi\!\left(\sum_{i=1}^{n-1}u_{i,i+1}+0\right)=\psi(u),\label{eq:psi-block}
\end{align}
where the second equality uses that the entry $[\mathrm{diag}(u,1)]_{n,n+1}=0$ by the block-diagonal structure (the row $n$ of $\mathrm{diag}(u,1)$ has $u_{n1},\ldots,u_{nn}$ in the first $n$ columns and $0$ in column $n+1$).

By~\eqref{eq:Whittaker-equivariance} (the Whittaker equivariance of $W$), \eqref{eq:diag-mult}, and~\eqref{eq:psi-block}:
\begin{align}\label{eq:W-shift}
W(\mathrm{diag}(ug,1)u_Q)&=W(\mathrm{diag}(u,1)\mathrm{diag}(g,1)u_Q)\\
&=\psi^{-1}(\mathrm{diag}(u,1))\,W(\mathrm{diag}(g,1)u_Q)\notag\\
&=\psi^{-1}(u)\,W(\mathrm{diag}(g,1)u_Q).\notag
\end{align}
Similarly $V(ug)=\psi(u)V(g)$ by~\eqref{eq:Whittaker-equivariance}. Finally $|\det(ug)|=|\det u|\cdot|\det g|=|\det g|$ since $u$ is unipotent. Multiplying:
$$W(\mathrm{diag}(ug,1)u_Q)V(ug)|\det(ug)|^{s-\tfrac12}=\psi^{-1}(u)\psi(u)\,W(\mathrm{diag}(g,1)u_Q)V(g)|\det g|^{s-\tfrac12}.$$
Since $\psi^{-1}(u)\psi(u)=1$, the integrand is left-$N_n$-invariant. \qed
\end{proof}

Equation~\eqref{eq:Whittaker-equivariance} referenced above is the standard Whittaker equivariance, which we make explicit:
\begin{equation}\label{eq:Whittaker-equivariance}
W(n_{n+1}\,h)=\psi^{-1}(n_{n+1})\,W(h),\quad V(n_n\,g)=\psi(n_n)\,V(g),
\end{equation}
for all $n_{n+1}\in N_{n+1},\,h\in\mathrm{GL}_{n+1}(F)$, $n_n\in N_n,\,g\in\mathrm{GL}_n(F)$, valid by definition of the $\psi^{-1}$-Whittaker model of $\Pi$ and the $\psi$-Whittaker model of $\pi$.

\subsection{Existence of newforms in the required Whittaker models}

\begin{lemma}\label{lem:newform-exists}
Both $W_\Pi^{\mathrm{new}}\in\mathcal{W}(\Pi,\psi^{-1})$ and $V_\pi^{\mathrm{new}}\in\mathcal{W}(\pi,\psi)$ exist, are unique up to scalar in their respective Whittaker models, and have nonzero value at the identity matrix.
\end{lemma}

\begin{proof}
We invoke Theorem~\ref{thm:JPSS-newform} twice:
\begin{itemize}
\item Apply with $(\sigma,\psi)\leftarrow(\Pi,\psi^{-1})$. The character $\psi^{-1}$ is nontrivial (since $\psi$ is) with conductor $\mathfrak{o}$ (the conductor of $\psi^{-1}$ equals that of $\psi$ because $\psi^{-1}(x)=\overline{\psi(x)}$ for $\mathbb{C}$-valued $\psi$, and conjugation preserves the conductor exponent). $\Pi$ is generic irreducible admissible by hypothesis. Theorem~\ref{thm:JPSS-newform} gives the conductor exponent $c(\Pi)$ and a one-dimensional space of $K_{n+1}(\mathfrak{p}^{c(\Pi)})$-fixed Whittaker functions; part (iv) gives nonvanishing at $I_{n+1}$.
\item Apply with $(\sigma,\psi)\leftarrow(\pi,\psi)$. $\pi$ is generic irreducible admissible by hypothesis, $\psi$ has conductor $\mathfrak{o}$. The same argument yields $V_\pi^{\mathrm{new}}$ unique up to scalar with $V_\pi^{\mathrm{new}}(I_n)\ne 0$.
\end{itemize}
After normalization $W_\Pi^{\mathrm{new}}(I_{n+1})=V_\pi^{\mathrm{new}}(I_n)=1$, both newforms are determined uniquely. \qed
\end{proof}

\subsection{Independence of the choice of generator $Q$}

\begin{remark}\label{rem:uQ-Q-class}
The conclusion of Theorem~\ref{thm:Matringe} (i.e., the identity $\Psi_{u_Q}=L$) is stated for \emph{any} generator $Q$ of $\mathfrak{q}^{-1}$. Different choices of $Q$ (differing by a unit in $\mathfrak{o}^\times$) thus lead to the same value of the right-hand side $L(s,\Pi\times\pi)$, which is intrinsic to $(\Pi,\pi)$ and does not depend on $Q$. By Matringe's theorem the integrals $\Psi_{u_Q}(s,W_\Pi^{\mathrm{new}},V_\pi^{\mathrm{new}})$ for varying generators $Q$ all equal $L(s,\Pi\times\pi)$, hence equal each other. (This compatibility is part of the content of Matringe's theorem; we do not need a separate proof.)
\end{remark}

\subsection{Surjectivity of the parametrization $s\mapsto q^{-s}$}

\begin{lemma}\label{lem:q-exp-surjective}
The map $\mathbb{C}\to\mathbb{C}^\times$, $s\mapsto q^{-s}=\exp(-s\log q)$, is surjective.
\end{lemma}

\begin{proof}
$\log q\in\mathbb{R}_{>0}$ since $q\ge 2$. The complex exponential $z\mapsto e^z$ is surjective onto $\mathbb{C}^\times$: given $w=re^{i\theta}\in\mathbb{C}^\times$ with $r>0,\theta\in\mathbb{R}$, set $z=\log r+i\theta\in\mathbb{C}$; then $e^z=w$. Composing with the bijection $s\mapsto -s\log q$ on $\mathbb{C}$ shows $s\mapsto q^{-s}$ is surjective onto $\mathbb{C}^\times$. \qed
\end{proof}

\section{Robustness: Boundary and Degenerate Cases}

We confirm that the proof of Theorem~\ref{thm:main} covers all generic irreducible admissible $\pi$ on $\mathrm{GL}_n(F)$, including degenerate-looking cases.

\subsection{Unramified $\pi$: $c(\pi)=0$}

If $\pi$ is unramified, then $c(\pi)=0$ and $\mathfrak{q}=\mathfrak{q}^{-1}=\mathfrak{o}$. A generator of $\mathfrak{o}$ as a fractional ideal is any unit; the canonical choice is $Q=1$, giving $u_Q=I_{n+1}+E_{n,n+1}$. Remark~\ref{rem:uQ-Q-class} ensures the conclusion is independent of the choice of unit $Q$. The newform $V_\pi^{\mathrm{new}}$ is the spherical Whittaker function of $\pi$ (i.e., the $\mathrm{GL}_n(\mathfrak{o})$-fixed vector). Theorem~\ref{thm:Matringe} applies and the conclusion holds. (When $\Pi$ is also unramified, this case recovers the classical Casselman--Shalika identity.)

\subsection{Highly ramified $\pi$: $c(\pi)$ large}

For $c(\pi)\gg 0$, the generator $Q$ of $\mathfrak{q}^{-1}=\mathfrak{p}^{-c(\pi)}$ has large absolute value $|Q|=q^{c(\pi)}$. Theorem~\ref{thm:Matringe} is stated uniformly in $c(\pi)\ge 0$, so the conclusion is unaffected by the size of the conductor; $\Psi(s,W,V)=L(s,\Pi\times\pi)$ continues to hold.

\subsection{Supercuspidal, discrete-series, principal series, and induced $\pi$}

The classification of generic irreducible admissible representations of $\mathrm{GL}_n(F)$ (Bernstein--Zelevinsky) includes:
\begin{itemize}
\item supercuspidal representations,
\item discrete series (Steinberg-type),
\item parabolically induced representations from generic representations of Levi subgroups.
\end{itemize}
Theorem~\ref{thm:Matringe} places no restriction on the type of $\pi$; the identity $\Psi=L$ holds in all cases.

\subsection{Edge cases of dimensions}

For $n=1$, $\mathrm{GL}_1(F)=F^\times$ and a generic irreducible admissible representation $\pi$ is simply a quasi-character $\chi:F^\times\to\mathbb{C}^\times$. The Whittaker model is the one-dimensional space $\mathbb{C}\cdot\chi$, and the integral~\eqref{eq:RS-integral} becomes
$$\Psi(s,W,V)=\int_{F^\times}W(\mathrm{diag}(g,1)u_Q)\,\chi(g)\,|g|^{s-\tfrac12}\,d^\times g.$$
Theorem~\ref{thm:Matringe} applies; the result equals $L(s,\Pi\times\chi)$, which has the form~\eqref{eq:L-shape-final} and is hence finite and nonzero in the sense of Definition~\ref{def:finite-nonzero}.

The case $n=0$ is vacuous: $\mathrm{GL}_0(F)$ is the trivial group, and the only admissible representation is the trivial one of trivial dimension. The problem implicitly assumes $n\ge 1$.

\section{Why the Newform is the Natural Choice}

\begin{remark}\label{rem:why-newform}
For an arbitrary $W\in\mathcal{W}(\Pi,\psi^{-1})$ and $V\in\mathcal{W}(\pi,\psi)$, Theorem~\ref{thm:JPSS-RS}(iii) only guarantees that $\Psi(s,W,V)\in\mathbb{C}[q^{-s},q^s]\cdot L(s,\Pi\times\pi)$. Concretely, for some Laurent polynomial $Q_{W,V}\in\mathbb{C}[X,X^{-1}]$,
$$\Psi(s,W,V)=Q_{W,V}(q^{-s})\cdot L(s,\Pi\times\pi).$$
If $Q_{W,V}$ has any zero in $\mathbb{C}^\times$, the integral $\Psi(s,W,V)$ vanishes at the corresponding values of $s$, violating Definition~\ref{def:finite-nonzero}(a). The point of Matringe's theorem is precisely that $(W_\Pi^{\mathrm{new}},V_\pi^{\mathrm{new}})$, paired with the conductor twist $u_Q$, gives $Q_{W,V}\equiv 1$, hence $\Psi=L$ exactly. The conductor twist $u_Q$ is essential: with $u_Q$ replaced by $I_{n+1}$, the integral $\Psi(s,W_\Pi^{\mathrm{new}},V_\pi^{\mathrm{new}})$ is generally a proper polynomial multiple of $L(s,\Pi\times\pi)$, with possible zeros.
\end{remark}

\section{Conclusion}

We have shown:
\begin{enumerate}
\item[(1)] An explicit witness $W=W_\Pi^{\mathrm{new}}\in\mathcal{W}(\Pi,\psi^{-1})$ satisfies the required property in Theorem~\ref{thm:main}.
\item[(2)] For every generic irreducible admissible $\pi$ on $\mathrm{GL}_n(F)$, the choice $V=V_\pi^{\mathrm{new}}$ gives, by Matringe's identity~\eqref{eq:Matringe-identity}, the equality
$$\Psi(s,W,V)\;=\;L(s,\Pi\times\pi).$$
\item[(3)] $L(s,\Pi\times\pi)$ is a nonzero rational function in $q^{-s}$ of the form $1/P(q^{-s})$ with $P(0)=1$, hence has no zeros in $\mathbb{C}$.
\end{enumerate}
Combined, these three facts establish that $\Psi(s,W,V)$ is finite and nonzero for all $s\in\mathbb{C}$ in the precise sense of Definition~\ref{def:finite-nonzero}, answering the problem affirmatively.

\medskip

\noindent\textbf{References.}
\begin{itemize}
\item H.\ Jacquet, I.\ I.\ Piatetski-Shapiro, J.\ A.\ Shalika, \emph{Conducteur des repr\'esentations du groupe lin\'eaire}, Math.\ Ann.\ \textbf{256} (1981), 199--214.
\item H.\ Jacquet, I.\ I.\ Piatetski-Shapiro, J.\ A.\ Shalika, \emph{Rankin-Selberg convolutions}, Amer.\ J.\ Math.\ \textbf{105} (1983), 367--464.
\item H.\ Jacquet, \emph{A correction to ``Conducteur des repr\'esentations du groupe lin\'eaire''}, Pacific J.\ Math.\ \textbf{260} (2012), no.\ 2, 515--525.
\item N.\ Matringe, \emph{Essential Whittaker functions for $\mathrm{GL}(n)$}, Doc.\ Math.\ \textbf{18} (2013), 1191--1214; arXiv:1201.5506.
\end{itemize}

\end{document}
